\documentclass[11pt,a4paper]{article}

%================================================
% PACKAGES
%================================================
\usepackage[utf8]{inputenc}
\usepackage[T1]{fontenc}
\usepackage[french]{babel}
\usepackage[a4paper,margin=1.9cm]{geometry}

\usepackage{amsmath,amssymb,amsthm,mathtools}
\usepackage{siunitx}
\usepackage{physics}
\usepackage{lmodern}
\usepackage{xcolor}
\usepackage{hyperref}
\usepackage{fancyhdr}
\usepackage{lastpage}
\usepackage{enumitem}
\usepackage{tcolorbox}
\usepackage{booktabs}
\usepackage{array}
\usepackage{multicol}
\usepackage{tikz}
\usepackage{pgfplots}
\usepackage{chemfig}
\usepackage{titlesec}
\usepackage{stmaryrd}

\pgfplotsset{compat=1.18}

%================================================
% MISE EN PAGE
%================================================
\pagestyle{fancy}
\fancyhf{}
\fancyhead[L]{Physique-Chimie -- Première Spécialité}
\fancyhead[R]{Feuille d'exercices corrigés}
\fancyfoot[C]{\thepage\ /\pageref{LastPage}}

\hypersetup{
    colorlinks=true,
    linkcolor=blue!60!black,
    urlcolor=blue!60!black,
    pdftitle={Exercices corrigés de physique-chimie - Première spécialité},
    pdfauthor={},
    pdfsubject={Exercices et problèmes avec corrigés intégrés}
}

\setlength{\parindent}{0pt}
\setlength{\parskip}{0.45em}

\titleformat{\section}{\Large\bfseries\color{blue!60!black}}{\thesection.}{0.6em}{}
\titleformat{\subsection}{\large\bfseries\color{blue!45!black}}{\thesubsection.}{0.6em}{}
\titleformat{\subsubsection}{\normalsize\bfseries\color{blue!30!black}}{\thesubsubsection.}{0.6em}{}

%================================================
% BOÎTES
%================================================
\newtcolorbox{objectifbox}{
    colback=blue!5,
    colframe=blue!60!black,
    title=Objectif,
    fonttitle=\bfseries
}

\newtcolorbox{enoncebox}{
    colback=gray!4,
    colframe=black!70,
    title=Énoncé,
    fonttitle=\bfseries
}

\newtcolorbox{corrigebox}{
    colback=green!5,
    colframe=green!45!black,
    title=Corrigé,
    fonttitle=\bfseries
}

\newtcolorbox{methodebox}{
    colback=orange!8,
    colframe=orange!70!black,
    title=Méthode / idée de fond,
    fonttitle=\bfseries
}

\newtcolorbox{attentionbox}{
    colback=red!5,
    colframe=red!65!black,
    title=Point de vigilance,
    fonttitle=\bfseries
}

\newtcolorbox{bilanbox}{
    colback=cyan!6,
    colframe=cyan!55!black,
    title=Bilan,
    fonttitle=\bfseries
}

%================================================
% COMMANDES
%================================================
\newcommand{\Rlim}{réactif limitant}
\newcommand{\Hplus}{\mathrm{H^+}}
\newcommand{\OHmoins}{\mathrm{HO^-}}
\newcommand{\eau}{\mathrm{H_2O}}
\newcommand{\COdeux}{\mathrm{CO_2}}
\newcommand{\Odeux}{\mathrm{O_2}}
\newcommand{\Hdeux}{\mathrm{H_2}}
\newcommand{\NadeuxSOquatre}{\mathrm{Na_2SO_4}}

%================================================
% DOCUMENT
%================================================
\begin{document}

\begin{center}
    {\LARGE \textbf{Grande feuille d'exercices corrigés}}\\[0.2em]
    {\Large \textbf{Physique-Chimie -- Première Spécialité}}\\[0.2em]
    {\normalsize Plusieurs chapitres -- raisonnements, problèmes, compréhension fine}\\[0.6em]
    \rule{0.9\linewidth}{0.6pt}
\end{center}

\begin{objectifbox}
Ce document a été conçu pour :
\begin{itemize}
    \item faire travailler plusieurs chapitres importants de la spécialité ;
    \item tester la compréhension réelle des notions, pas seulement l'application d'une recette ;
    \item proposer des exercices variés : questions de cours, calculs, interprétations, problèmes ouverts ;
    \item placer le corrigé juste après chaque exercice pour que l'élève puisse comparer son raisonnement.
\end{itemize}
\end{objectifbox}

\tableofcontents
\vspace{1em}

%================================================
\section{Chapitre 1 -- Structure de la matière, polarité, interactions}
%================================================

%------------------------------------------------
\subsection{Exercice 1 -- Liaisons polaires ou molécule polaire ?}
%------------------------------------------------

\begin{enoncebox}
On considère les molécules suivantes :
\[
\mathrm{H_2O},\qquad \mathrm{CO_2},\qquad \mathrm{CH_4},\qquad \mathrm{NH_3}
\]

On admet les géométries suivantes :
\begin{itemize}
    \item \(\mathrm{H_2O}\) : coudée ;
    \item \(\mathrm{CO_2}\) : linéaire ;
    \item \(\mathrm{CH_4}\) : tétraédrique symétrique ;
    \item \(\mathrm{NH_3}\) : pyramidale.
\end{itemize}

\begin{enumerate}[label=\arabic*)]
    \item Dire quelles liaisons sont polaires.
    \item Dire quelles molécules sont polaires.
    \item Expliquer pourquoi une molécule peut contenir des liaisons polaires tout en étant apolaire.
    \item Comparer le comportement attendu de \(\mathrm{H_2O}\) et \(\mathrm{CO_2}\) vis-à-vis de l'eau.
\end{enumerate}
\end{enoncebox}

\begin{corrigebox}
\textbf{1) Liaisons polaires}

Une liaison est polaire lorsqu'il existe une différence notable d'électronégativité entre les deux atomes liés.
\begin{itemize}
    \item Dans \(\mathrm{H_2O}\), les liaisons O--H sont polaires.
    \item Dans \(\mathrm{CO_2}\), les liaisons C=O sont polaires.
    \item Dans \(\mathrm{CH_4}\), les liaisons C--H sont très peu polaires ; à ce niveau on les considère souvent comme non polaires ou très faiblement polaires.
    \item Dans \(\mathrm{NH_3}\), les liaisons N--H sont polaires.
\end{itemize}

\textbf{2) Molécules polaires}

\begin{itemize}
    \item \(\mathrm{H_2O}\) est polaire : géométrie coudée, les effets des deux liaisons ne se compensent pas.
    \item \(\mathrm{CO_2}\) est apolaire : les deux liaisons C=O sont opposées et se compensent.
    \item \(\mathrm{CH_4}\) est apolaire : géométrie symétrique.
    \item \(\mathrm{NH_3}\) est polaire : géométrie pyramidale, pas de compensation.
\end{itemize}

\textbf{3) Pourquoi des liaisons polaires ne suffisent-elles pas ?}

La polarité globale dépend :
\begin{itemize}
    \item de la polarité de chaque liaison ;
    \item de la géométrie de la molécule.
\end{itemize}
Si les effets des liaisons se compensent vectoriellement, la molécule peut être globalement apolaire.

\textbf{4) Comparaison \(\mathrm{H_2O}\) / \(\mathrm{CO_2}\)}

L'eau est polaire et peut former des interactions fortes avec d'autres espèces polaires ou ioniques.  
Le dioxyde de carbone est apolaire dans son ensemble ; ses interactions avec l'eau sont moins favorables.  
On s'attend donc à ce que l'eau interagisse bien plus fortement avec elle-même qu'avec \(\mathrm{CO_2}\).

\begin{bilanbox}
La phrase clé à retenir est :
\[
\text{Liaison polaire} \neq \text{forcément molécule polaire}.
\]
\end{bilanbox}
\end{corrigebox}

%------------------------------------------------
\subsection{Exercice 2 -- Solubilité et raisonnement microscopique}
%------------------------------------------------

\begin{enoncebox}
On souhaite prévoir qualitativement la solubilité dans l'eau des espèces suivantes :
\[
\mathrm{NaCl},\qquad \mathrm{I_2},\qquad \text{éthanol},\qquad \text{huile}
\]

On rappelle que :
\begin{itemize}
    \item l'eau est un solvant polaire ;
    \item l'éthanol possède un groupe O--H ;
    \item l'huile est constituée majoritairement d'espèces apolaires.
\end{itemize}

\begin{enumerate}[label=\arabic*)]
    \item Classer ces espèces selon leur aptitude à se mélanger à l'eau.
    \item Justifier chaque réponse en termes d'interactions microscopiques.
    \item Expliquer pourquoi la règle ``le semblable dissout le semblable'' ne doit pas être récitée sans réflexion.
\end{enumerate}
\end{enoncebox}

\begin{corrigebox}
\textbf{1) Classement qualitatif}

\begin{itemize}
    \item \(\mathrm{NaCl}\) : très soluble.
    \item Éthanol : très soluble ou miscible.
    \item \(\mathrm{I_2}\) : peu soluble.
    \item Huile : non miscible ou très peu miscible.
\end{itemize}

\textbf{2) Justification}

\begin{itemize}
    \item \(\mathrm{NaCl}\) : solide ionique ; l'eau peut hydrater les ions \(\mathrm{Na^+}\) et \(\mathrm{Cl^-}\).
    \item Éthanol : la présence du groupe O--H permet des interactions favorables avec l'eau, notamment des liaisons hydrogène.
    \item \(\mathrm{I_2}\) : molécule apolaire ; les interactions avec l'eau sont peu favorables.
    \item Huile : espèces majoritairement apolaires, donc très mauvaises interactions avec l'eau.
\end{itemize}

\textbf{3) Pourquoi la règle est-elle insuffisante seule ?}

La règle est utile comme idée directrice, mais elle ne vaut que si on la traduit physiquement :
\begin{itemize}
    \item quelles interactions faut-il rompre ?
    \item quelles interactions nouvelles se forment ?
    \item ces interactions nouvelles compensent-elles les anciennes ?
\end{itemize}

\begin{methodebox}
Comprendre la dissolution, ce n'est pas apprendre une phrase par coeur ; c'est comparer des interactions.
\end{methodebox}
\end{corrigebox}

%------------------------------------------------
\subsection{Problème 1 -- Pourquoi l'eau et l'huile se séparent-elles ?}
%------------------------------------------------

\begin{enoncebox}
On mélange vigoureusement de l'eau et de l'huile dans un flacon transparent.  
Pendant quelques secondes, on observe une dispersion trouble, puis deux phases distinctes réapparaissent.

\begin{enumerate}[label=\arabic*)]
    \item Expliquer pourquoi l'agitation permet temporairement de former une dispersion.
    \item Expliquer pourquoi cette dispersion n'est pas stable.
    \item Montrer comment cette expérience mobilise simultanément les notions de polarité, cohésion et miscibilité.
\end{enumerate}
\end{enoncebox}

\begin{corrigebox}
\textbf{a) Effet de l'agitation}

L'agitation apporte de l'énergie mécanique et fragmente un liquide en fines gouttelettes dispersées dans l'autre.  
On force donc provisoirement un état plus mélangé.

\textbf{b) Pourquoi la dispersion disparaît-elle ?}

Les interactions eau--eau sont fortes car l'eau est polaire et peut former des liaisons hydrogène.  
Les interactions huile--huile sont plus favorables que les interactions eau--huile.  
Le système revient donc vers l'état où chaque liquide interagit préférentiellement avec lui-même.

\textbf{c) Lien avec les notions du cours}

\begin{itemize}
    \item \textbf{Polarité} : eau polaire, huile apolaire.
    \item \textbf{Cohésion} : chaque liquide conserve une cohésion interne.
    \item \textbf{Miscibilité} : les interactions croisées eau--huile ne sont pas assez favorables pour stabiliser durablement un mélange homogène.
\end{itemize}

\begin{bilanbox}
Un mélange homogène ne dépend pas seulement de l'action de mélanger ; il dépend surtout des interactions réellement possibles entre les espèces.
\end{bilanbox}
\end{corrigebox}

%================================================
\section{Chapitre 2 -- Transformations chimiques, quantité de matière, réactif limitant}
%================================================

%------------------------------------------------
\subsection{Exercice 3 -- Identifier le réactif limitant}
%------------------------------------------------

\begin{enoncebox}
On fait réagir du fer avec du dioxygène pour former de l'oxyde de fer(III), selon l'équation :
\[
4\,\mathrm{Fe} + 3\,\mathrm{O_2} \rightarrow 2\,\mathrm{Fe_2O_3}
\]

On dispose initialement de :
\[
n(\mathrm{Fe}) = \SI{0.80}{mol}, \qquad n(\mathrm{O_2}) = \SI{0.45}{mol}
\]

\begin{enumerate}[label=\arabic*)]
    \item Déterminer le réactif limitant.
    \item Calculer l'avancement maximal.
    \item Déterminer la quantité de matière finale de \(\mathrm{Fe_2O_3}\).
    \item Calculer la quantité de matière du réactif en excès restant à la fin.
\end{enumerate}
\end{enoncebox}

\begin{corrigebox}
On note \(x\) l'avancement.

\[
4\,\mathrm{Fe} + 3\,\mathrm{O_2} \rightarrow 2\,\mathrm{Fe_2O_3}
\]

\textbf{1) Réactif limitant}

Les contraintes sont :
\[
x \leq \frac{0.80}{4} = 0.20
\]
et
\[
x \leq \frac{0.45}{3} = 0.15
\]

La plus petite valeur est \(\SI{0.15}{mol}\).  
Donc :
\[
x_{\max} = \SI{0.15}{mol}
\]
Le dioxygène est le \Rlim.

\textbf{2) Avancement maximal}

\[
x_{\max} = \SI{0.15}{mol}
\]

\textbf{3) Quantité de matière finale de \(\mathrm{Fe_2O_3}\)}

Le coefficient stoechiométrique du produit vaut 2, donc :
\[
n_f(\mathrm{Fe_2O_3}) = 2x_{\max} = 2\times 0.15 = \SI{0.30}{mol}
\]

\textbf{4) Réactif en excès restant}

Pour le fer :
\[
n_f(\mathrm{Fe}) = 0.80 - 4x_{\max} = 0.80 - 4\times 0.15 = 0.80 - 0.60 = \SI{0.20}{mol}
\]

\begin{attentionbox}
L'erreur classique consiste à comparer directement \(0.80\) et \(0.45\), ce qui n'a pas de sens sans tenir compte des coefficients de l'équation-bilan.
\end{attentionbox}
\end{corrigebox}

%------------------------------------------------
\subsection{Exercice 4 -- Raisonner sans tableau automatique}
%------------------------------------------------

\begin{enoncebox}
Le méthane brûle dans le dioxygène selon :
\[
\mathrm{CH_4} + 2\,\mathrm{O_2} \rightarrow \mathrm{CO_2} + 2\,\mathrm{H_2O}
\]

On dispose de :
\[
n(\mathrm{CH_4}) = \SI{0.50}{mol}, \qquad n(\mathrm{O_2}) = \SI{1.50}{mol}
\]

Sans faire de tableau d'avancement complet au départ :
\begin{enumerate}[label=\arabic*)]
    \item dire si le méthane peut être entièrement consommé ;
    \item préciser quel réactif est en excès ;
    \item calculer les quantités de matière finales des produits ;
    \item vérifier ensuite la cohérence du raisonnement.
\end{enumerate}
\end{enoncebox}

\begin{corrigebox}
\textbf{1) Le méthane peut-il être entièrement consommé ?}

Pour consommer \(\SI{0.50}{mol}\) de méthane, il faut :
\[
n(\mathrm{O_2})_{\text{nécessaire}} = 2\times 0.50 = \SI{1.00}{mol}
\]

Or on dispose de \(\SI{1.50}{mol}\) de dioxygène.  
Donc le méthane peut être totalement consommé.

\textbf{2) Réactif en excès}

Le dioxygène est en excès.

\textbf{3) Quantités finales des produits}

Si tout le méthane réagit :
\[
n_f(\mathrm{CO_2}) = \SI{0.50}{mol}
\]
car le coefficient stoechiométrique est 1.

Et :
\[
n_f(\mathrm{H_2O}) = 2\times 0.50 = \SI{1.00}{mol}
\]

Dioxygène restant :
\[
n_f(\mathrm{O_2}) = 1.50 - 1.00 = \SI{0.50}{mol}
\]

\textbf{4) Vérification}

Le raisonnement est cohérent car :
\begin{itemize}
    \item le méthane a bien disparu ;
    \item le dioxygène reste en quantité positive ;
    \item les quantités de produits suivent l'équation-bilan.
\end{itemize}

\begin{methodebox}
Savoir faire un tableau est utile.  
Mais comprendre avant tout les proportions stoechiométriques est encore plus important.
\end{methodebox}
\end{corrigebox}

%------------------------------------------------
\subsection{Problème 2 -- Une réaction n'est pas juste une formule}
%------------------------------------------------

\begin{enoncebox}
Dans un laboratoire, on souhaite produire du dihydrogène par réaction entre un métal et un acide.  
On fait réagir du zinc avec une solution aqueuse acide selon :
\[
\mathrm{Zn} + 2\,\mathrm{H^+} \rightarrow \mathrm{Zn^{2+}} + \mathrm{H_2}
\]

On introduit :
\[
n(\mathrm{Zn}) = \SI{0.12}{mol},\qquad n(\mathrm{H^+}) = \SI{0.18}{mol}
\]

\begin{enumerate}[label=\arabic*)]
    \item Déterminer le réactif limitant.
    \item Calculer la quantité de matière de dihydrogène formée.
    \item Calculer la quantité de matière de l'espèce en excès restante.
    \item Expliquer pourquoi cette réaction permet de transformer une espèce ionique dissoute en gaz observable.
\end{enumerate}
\end{enoncebox}

\begin{corrigebox}
\textbf{1) Réactif limitant}

\[
\mathrm{Zn} + 2\,\mathrm{H^+} \rightarrow \mathrm{Zn^{2+}} + \mathrm{H_2}
\]

Si tout le zinc réagit, il faut :
\[
n(\Hplus)_{\text{nécessaire}} = 2\times 0.12 = \SI{0.24}{mol}
\]
Or on ne dispose que de \(\SI{0.18}{mol}\) de \(\Hplus\).  
Donc \(\Hplus\) est le \Rlim.

\textbf{2) Quantité de matière de dihydrogène formée}

Pour 2 mol de \(\Hplus\), on forme 1 mol de \(\Hdeux\).  
Donc :
\[
n(\Hdeux) = \frac{0.18}{2} = \SI{0.09}{mol}
\]

\textbf{3) Zinc restant}

Le zinc consommé vaut :
\[
n(\mathrm{Zn})_{\text{consommé}} = \SI{0.09}{mol}
\]
car le coefficient du zinc est 1 pour 1 mol de \(\Hdeux\).

Donc :
\[
n_f(\mathrm{Zn}) = 0.12 - 0.09 = \SI{0.03}{mol}
\]

\textbf{4) Sens physique de la réaction}

Cette transformation convertit des ions \(\Hplus\), invisibles en solution, en molécules de dihydrogène \(\Hdeux\), qui peuvent former des bulles et devenir observables à notre échelle.  
La réaction illustre donc comment une évolution microscopique se traduit par un effet macroscopique visible.

\begin{bilanbox}
Une équation chimique ne sert pas seulement à calculer : elle raconte une transformation réelle de la matière.
\end{bilanbox}
\end{corrigebox}

%================================================
\section{Chapitre 3 -- Solutions, concentration, dilution}
%================================================

%------------------------------------------------
\subsection{Exercice 5 -- Préparer une solution diluée}
%------------------------------------------------

\begin{enoncebox}
On dispose d'une solution mère de sulfate de cuivre de concentration :
\[
C_m = \SI{0.50}{mol.L^{-1}}
\]

On veut préparer un volume :
\[
V_f = \SI{100.0}{mL}
\]
de solution fille de concentration :
\[
C_f = \SI{0.080}{mol.L^{-1}}
\]

\begin{enumerate}[label=\arabic*)]
    \item Calculer le volume de solution mère à prélever.
    \item Décrire la préparation expérimentale.
    \item Expliquer pourquoi la quantité de soluté prélevée est conservée lors de la dilution.
\end{enumerate}
\end{enoncebox}

\begin{corrigebox}
\textbf{1) Calcul du volume prélevé}

On utilise :
\[
C_m V_p = C_f V_f
\]
avec \(V_p\) le volume prélevé.

\[
V_p = \frac{C_f V_f}{C_m}
\]

En prenant \(V_f = \SI{0.1000}{L}\) :
\[
V_p = \frac{0.080 \times 0.1000}{0.50} = \SI{0.0160}{L} = \SI{16.0}{mL}
\]

\textbf{2) Préparation}

\begin{itemize}
    \item Prélever \(\SI{16.0}{mL}\) de solution mère à l'aide d'une pipette graduée ou d'une pipette jaugée adaptée.
    \item Introduire ce volume dans une fiole jaugée de \(\SI{100.0}{mL}\).
    \item Compléter avec de l'eau distillée jusqu'au trait de jauge.
    \item Boucher, puis homogénéiser.
\end{itemize}

\textbf{3) Pourquoi la quantité de soluté est-elle conservée ?}

Lors d'une dilution, on ajoute du solvant mais on ne retire ni ne crée de soluté.  
La quantité de matière de soluté dissous prélevée au départ est donc la même dans la solution fille finale.

\begin{attentionbox}
Une dilution diminue la concentration, pas la quantité de matière de soluté présente dans le volume transféré.
\end{attentionbox}
\end{corrigebox}

%------------------------------------------------
\subsection{Exercice 6 -- Dissolution ou dilution ?}
%------------------------------------------------

\begin{enoncebox}
Pour chacune des situations suivantes, préciser s'il s'agit d'une dissolution ou d'une dilution, puis expliquer la différence de sens physique.

\begin{enumerate}[label=\alph*)]
    \item On introduit \(\SI{5.0}{g}\) de chlorure de sodium solide dans de l'eau.
    \item On prélève \(\SI{10.0}{mL}\) d'une solution de glucose concentrée puis on ajoute de l'eau jusqu'à \(\SI{100.0}{mL}\).
    \item On met un morceau de sucre dans un café chaud.
\end{enumerate}

Puis répondre :
\begin{enumerate}[label=\arabic*)]
    \item Quelle grandeur est conservée lors d'une dilution ?
    \item Quelle grandeur augmente lors d'une dilution ?
    \item Pourquoi les deux opérations ne doivent-elles pas être confondues ?
\end{enumerate}
\end{enoncebox}

\begin{corrigebox}
\textbf{Identification}

\begin{itemize}
    \item \textbf{a)} Dissolution : on part d'un soluté solide et on le disperse dans un solvant.
    \item \textbf{b)} Dilution : on part déjà d'une solution et on ajoute du solvant pour diminuer sa concentration.
    \item \textbf{c)} Dissolution : le sucre solide se dissout dans le café.
\end{itemize}

\textbf{1) Grandeur conservée lors d'une dilution}

La quantité de matière de soluté contenue dans le prélèvement est conservée.

\textbf{2) Grandeur qui augmente lors d'une dilution}

Le volume total de solution augmente.

\textbf{3) Pourquoi ne pas confondre ?}

Dans une dissolution, on introduit du soluté dans un solvant.  
Dans une dilution, on possède déjà une solution et on change seulement sa concentration en ajoutant du solvant.

\begin{bilanbox}
\[
\text{Dissolution} \neq \text{Dilution}
\]
La première crée la solution à partir d'un soluté ; la seconde modifie une solution déjà préparée.
\end{bilanbox}
\end{corrigebox}

%================================================
\section{Chapitre 4 -- Électricité, puissance, énergie}
%================================================

%------------------------------------------------
\subsection{Exercice 7 -- Comprendre la puissance électrique}
%------------------------------------------------

\begin{enoncebox}
Un appareil électrique fonctionne sous une tension de \(\SI{230}{V}\) et est parcouru par un courant d'intensité \(\SI{3.0}{A}\).

\begin{enumerate}[label=\arabic*)]
    \item Calculer sa puissance électrique.
    \item Calculer l'énergie électrique reçue en \(\SI{15}{min}\).
    \item Expliquer la différence entre puissance et énergie.
    \item Un élève dit : ``plus la puissance est grande, plus l'énergie est grande''. Expliquer pourquoi cette phrase est incomplète.
\end{enumerate}
\end{enoncebox}

\begin{corrigebox}
\textbf{1) Puissance}

\[
P = UI = 230 \times 3.0 = \SI{690}{W}
\]

\textbf{2) Énergie reçue en \(\SI{15}{min}\)}

\[
E = P\Delta t
\]
avec :
\[
\Delta t = \SI{15}{min} = 15\times 60 = \SI{900}{s}
\]

Donc :
\[
E = 690 \times 900 = \SI{621000}{J}
\]

\textbf{3) Différence entre puissance et énergie}

\begin{itemize}
    \item La puissance mesure le débit de transfert d'énergie.
    \item L'énergie mesure la quantité totale transférée sur une durée.
\end{itemize}

\textbf{4) Pourquoi la phrase est-elle incomplète ?}

À durée fixée, une puissance plus grande implique une énergie plus grande.  
Mais l'énergie dépend aussi du temps de fonctionnement.  
Un appareil peu puissant qui fonctionne longtemps peut transférer plus d'énergie qu'un appareil très puissant utilisé brièvement.

\begin{methodebox}
La bonne phrase est :
\[
E = P\times \Delta t
\]
Il faut donc toujours penser simultanément à la puissance et à la durée.
\end{methodebox}
\end{corrigebox}

%------------------------------------------------
\subsection{Exercice 8 -- Batterie et autonomie}
%------------------------------------------------

\begin{enoncebox}
Une batterie d'énergie disponible
\[
E_{\text{bat}} = \SI{540}{kJ}
\]
alimente un dispositif électrique de puissance moyenne
\[
P = \SI{45}{W}
\]

\begin{enumerate}[label=\arabic*)]
    \item Estimer la durée de fonctionnement.
    \item Exprimer le résultat en secondes puis en heures.
    \item Expliquer pourquoi, dans un vrai système, cette durée est souvent surestimée par ce calcul simple.
\end{enumerate}
\end{enoncebox}

\begin{corrigebox}
\textbf{1) Durée de fonctionnement}

On utilise :
\[
E = P\Delta t \quad \Rightarrow \quad \Delta t = \frac{E}{P}
\]

Avec :
\[
E = \SI{540}{kJ} = \SI{540000}{J}
\]

Donc :
\[
\Delta t = \frac{540000}{45} = \SI{12000}{s}
\]

\textbf{2) En heures}

\[
\Delta t = \frac{12000}{3600} \approx \SI{3.33}{h}
\]

Soit environ :
\[
\SI{3}{h}\;20\;\text{min}
\]

\textbf{3) Pourquoi ce calcul surestime-t-il souvent ?}

Dans un système réel :
\begin{itemize}
    \item toute l'énergie stockée n'est pas toujours récupérable ;
    \item il existe des pertes ;
    \item la puissance n'est pas forcément constante ;
    \item le rendement n'est pas parfait.
\end{itemize}

\begin{attentionbox}
Les calculs idéaux sont indispensables pour raisonner, mais un système réel comporte presque toujours des pertes.
\end{attentionbox}
\end{corrigebox}

%------------------------------------------------
\subsection{Problème 3 -- Dimensionner un usage quotidien}
%------------------------------------------------

\begin{enoncebox}
Un élève souhaite alimenter une petite lampe LED pour lire le soir en camping.  
La lampe fonctionne avec une puissance moyenne de \(\SI{6.0}{W}\) et doit fonctionner \(\SI{2.5}{h}\) chaque soir pendant 4 jours.

On dispose de deux batteries :
\begin{itemize}
    \item batterie A : énergie disponible \(\SI{180}{kJ}\) ;
    \item batterie B : énergie disponible \(\SI{260}{kJ}\).
\end{itemize}

\begin{enumerate}[label=\arabic*)]
    \item Calculer l'énergie totale nécessaire.
    \item Dire quelle batterie choisir.
    \item Expliquer pourquoi un léger écart entre énergie calculée et énergie disponible peut être problématique en pratique.
\end{enumerate}
\end{enoncebox}

\begin{corrigebox}
\textbf{a) Énergie nécessaire}

Durée totale :
\[
\Delta t = 4\times 2.5 = \SI{10}{h}
\]

En secondes :
\[
\Delta t = 10\times 3600 = \SI{36000}{s}
\]

Énergie requise :
\[
E = P\Delta t = 6.0\times 36000 = \SI{216000}{J} = \SI{216}{kJ}
\]

\textbf{b) Choix de la batterie}

\begin{itemize}
    \item Batterie A : \(\SI{180}{kJ}\), insuffisante.
    \item Batterie B : \(\SI{260}{kJ}\), suffisante.
\end{itemize}

On choisit la batterie B.

\textbf{c) Pourquoi un faible écart peut-il poser problème ?}

Parce qu'en pratique :
\begin{itemize}
    \item le rendement n'est pas parfait ;
    \item la puissance peut varier ;
    \item une batterie ne délivre pas toujours exactement son énergie nominale ;
    \item il est prudent de garder une marge.
\end{itemize}

\begin{bilanbox}
Dimensionner un système, ce n'est pas seulement faire un calcul exact : c'est aussi prévoir les écarts du réel.
\end{bilanbox}
\end{corrigebox}

%================================================
\section{Chapitre 5 -- Mouvement, forces, énergie mécanique}
%================================================

%------------------------------------------------
\subsection{Exercice 9 -- Identifier les forces}
%------------------------------------------------

\begin{enoncebox}
Une caisse est tirée à vitesse constante sur un sol horizontal par une corde.  
On modélise la caisse comme un point matériel.

\begin{enumerate}[label=\arabic*)]
    \item Faire l'inventaire des forces appliquées à la caisse.
    \item Dire lesquelles sont des forces de contact et lesquelles sont des forces à distance.
    \item Expliquer pourquoi la vitesse constante donne une information importante sur le bilan des forces.
    \item Peut-on conclure que les forces qui s'exercent sont nulles ? Justifier.
\end{enumerate}
\end{enoncebox}

\begin{corrigebox}
\textbf{a) Forces appliquées}

\begin{itemize}
    \item le poids \(\vec{P}\), dû à l'attraction terrestre ;
    \item la réaction du sol ;
    \item la force de traction exercée par la corde ;
    \item éventuellement une force de frottement du sol sur la caisse.
\end{itemize}

\textbf{b) Contact ou distance}

\begin{itemize}
    \item Poids : force à distance.
    \item Réaction du sol : force de contact.
    \item Traction : force de contact.
    \item Frottement : force de contact.
\end{itemize}

\textbf{c) Rôle de la vitesse constante}

Si la vitesse est constante, l'accélération est nulle.  
On en déduit que la résultante des forces est nulle.

\textbf{d) Les forces sont-elles nulles ?}

Non.  
Les forces peuvent exister tout en se compensant.  
Vitesse constante signifie :
\[
\sum \vec{F} = \vec{0}
\]
et non pas que chaque force est nulle.

\begin{attentionbox}
Erreur classique : confondre ``résultante nulle'' et ``absence de forces''.
\end{attentionbox}
\end{corrigebox}

%------------------------------------------------
\subsection{Exercice 10 -- Énergie mécanique et frottements}
%------------------------------------------------

\begin{enoncebox}
Une bille descend une pente puis arrive sur une portion horizontale rugueuse.  
On observe qu'après la descente, sa vitesse diminue progressivement sur la partie horizontale jusqu'à s'annuler.

\begin{enumerate}[label=\arabic*)]
    \item Expliquer qualitativement l'évolution de l'énergie potentielle de pesanteur pendant la descente.
    \item Expliquer pourquoi la vitesse augmente d'abord.
    \item Expliquer ensuite pourquoi la vitesse diminue sur la partie horizontale.
    \item Dire si l'énergie mécanique du système bille-Terre se conserve sur toute la trajectoire.
\end{enumerate}
\end{enoncebox}

\begin{corrigebox}
\textbf{1) Pendant la descente}

L'altitude diminue, donc l'énergie potentielle de pesanteur diminue.

\textbf{2) Pourquoi la vitesse augmente-t-elle ?}

Une partie de l'énergie potentielle perdue se transforme en énergie cinétique.  
La bille accélère donc au cours de la descente.

\textbf{3) Pourquoi la vitesse diminue-t-elle ensuite ?}

Sur la partie horizontale rugueuse, les frottements dissipent de l'énergie.  
L'énergie cinétique diminue donc progressivement jusqu'à devenir nulle.

\textbf{4) Conservation de l'énergie mécanique ?}

Elle n'est pas conservée sur toute la trajectoire si des frottements interviennent.  
Sur la portion où les frottements sont négligeables, on peut considérer qu'elle se conserve approximativement.  
Sur la partie rugueuse, elle décroît.

\begin{methodebox}
Quand on parle d'énergie mécanique, la vraie question est toujours :
\begin{center}
\textit{les frottements sont-ils négligeables ou non ?}
\end{center}
\end{methodebox}
\end{corrigebox}

%------------------------------------------------
\subsection{Problème 4 -- Freinage d'un vélo}
%------------------------------------------------

\begin{enoncebox}
Un cycliste roule à vitesse non nulle sur une route horizontale.  
Il cesse de pédaler, puis freine progressivement jusqu'à l'arrêt.

\begin{enumerate}[label=\arabic*)]
    \item Décrire l'évolution de son énergie cinétique.
    \item Expliquer où ``va'' cette énergie.
    \item Pourquoi est-il faux de dire que l'énergie a disparu ?
    \item En quoi cette situation permet-elle de comprendre l'idée de dissipation ?
\end{enumerate}
\end{enoncebox}

\begin{corrigebox}
\textbf{1) Évolution de l'énergie cinétique}

L'énergie cinétique diminue au cours du ralentissement, puis devient nulle lorsque le vélo s'arrête.

\textbf{2) Où va cette énergie ?}

Elle est transférée vers d'autres formes :
\begin{itemize}
    \item échauffement des freins ;
    \item échauffement des pneus et de la route ;
    \item agitation thermique de l'air ;
    \item éventuellement un peu de son.
\end{itemize}

\textbf{3) Pourquoi ne disparaît-elle pas ?}

L'énergie ne disparaît pas : elle est transformée et transférée.  
Elle devient simplement moins facilement récupérable pour produire à nouveau le mouvement initial.

\textbf{4) Compréhension de la dissipation}

La dissipation correspond au fait qu'une énergie mécanique organisée se transforme en formes plus dispersées, en particulier en énergie thermique.

\begin{bilanbox}
Dire qu'une énergie est dissipée ne signifie pas qu'elle est détruite ; cela signifie qu'elle est transférée sous une forme moins exploitable pour le mouvement étudié.
\end{bilanbox}
\end{corrigebox}

%================================================
\section{Chapitre 6 -- Ondes, signaux, lumière}
%================================================

%------------------------------------------------
\subsection{Exercice 11 -- Grandeurs caractéristiques d'une onde}
%------------------------------------------------

\begin{enoncebox}
Une onde progressive périodique se propage à la célérité
\[
v = \SI{340}{m.s^{-1}}
\]
et a pour fréquence
\[
f = \SI{680}{Hz}
\]

\begin{enumerate}[label=\arabic*)]
    \item Calculer sa période.
    \item Calculer sa longueur d'onde.
    \item Expliquer le sens physique de la longueur d'onde.
    \item Dire ce qui change si la fréquence double alors que la célérité reste la même.
\end{enumerate}
\end{enoncebox}

\begin{corrigebox}
\textbf{1) Période}

\[
T = \frac{1}{f} = \frac{1}{680} \approx \SI{1.47e-3}{s}
\]

\textbf{2) Longueur d'onde}

\[
\lambda = \frac{v}{f} = \frac{340}{680} = \SI{0.50}{m}
\]

\textbf{3) Sens physique}

La longueur d'onde est la distance parcourue par l'onde pendant une période.  
C'est aussi la distance entre deux points successifs vibrant en phase.

\textbf{4) Si la fréquence double}

Si \(f\) double et \(v\) reste constante, alors :
\[
\lambda = \frac{v}{f}
\]
donc la longueur d'onde est divisée par 2.

\begin{attentionbox}
À célérité fixée :
\[
\lambda \propto \frac{1}{f}
\]
Fréquence et longueur d'onde varient en sens inverse.
\end{attentionbox}
\end{corrigebox}

%------------------------------------------------
\subsection{Exercice 12 -- Couleur, absorption, perception}
%------------------------------------------------

\begin{enoncebox}
On éclaire un objet avec de la lumière blanche.  
L'objet paraît rouge.

\begin{enumerate}[label=\arabic*)]
    \item Que peut-on dire qualitativement des radiations qu'il diffuse ou transmet ?
    \item Que devient la partie non rouge de la lumière reçue ?
    \item Pourquoi un même objet peut-il paraître différent selon l'éclairage ?
    \item Expliquer en quoi cette question mobilise un raisonnement physique, et pas seulement une connaissance de vocabulaire.
\end{enumerate}
\end{enoncebox}

\begin{corrigebox}
\textbf{1) Radiations diffusées ou transmises}

Si l'objet paraît rouge, cela signifie qu'il renvoie préférentiellement vers l'oeil les radiations rouges.

\textbf{2) Partie non rouge}

Les autres radiations sont majoritairement absorbées ou beaucoup moins diffusées/transmises.

\textbf{3) Pourquoi l'apparence dépend-elle de l'éclairage ?}

Parce que l'objet ne peut renvoyer que les radiations qu'il reçoit.  
Si la lumière incidente ne contient pas de rouge, un objet rouge ne pourra pas apparaître rouge de la même façon.

\textbf{4) Pourquoi est-ce un vrai raisonnement physique ?}

Il faut articuler :
\begin{itemize}
    \item la composition spectrale de la lumière incidente ;
    \item l'interaction lumière-matière ;
    \item la lumière effectivement reçue par l'oeil.
\end{itemize}
Ce n'est donc pas une simple phrase à apprendre.

\begin{bilanbox}
La couleur perçue n'est pas une propriété isolée de l'objet ; elle dépend aussi de la lumière qui l'éclaire.
\end{bilanbox}
\end{corrigebox}

%------------------------------------------------
\subsection{Problème 5 -- Un signal périodique et son interprétation}
%------------------------------------------------

\begin{enoncebox}
On enregistre un signal périodique à l'oscilloscope.  
On observe que le motif élémentaire se répète toutes les \(\SI{4.0}{ms}\).

\begin{enumerate}[label=\arabic*)]
    \item Déterminer la période.
    \item Calculer la fréquence.
    \item Expliquer pourquoi augmenter la fréquence change la sensation de hauteur d'un son.
    \item Un élève affirme : ``plus l'amplitude est grande, plus le son est aigu''. Discuter cette affirmation.
\end{enumerate}
\end{enoncebox}

\begin{corrigebox}
\textbf{1) Période}

Le motif se répète toutes les \(\SI{4.0}{ms}\), donc :
\[
T = \SI{4.0e-3}{s}
\]

\textbf{2) Fréquence}

\[
f = \frac{1}{T} = \frac{1}{4.0\times 10^{-3}} = \SI{250}{Hz}
\]

\textbf{3) Hauteur du son}

La hauteur d'un son est liée à sa fréquence :
\begin{itemize}
    \item fréquence plus grande \(\Rightarrow\) son plus aigu ;
    \item fréquence plus petite \(\Rightarrow\) son plus grave.
\end{itemize}

\textbf{4) Discussion de l'affirmation}

L'affirmation est fausse.  
L'amplitude est liée à l'intensité perçue, donc au caractère plus fort ou plus faible du son.  
La hauteur dépend de la fréquence, pas de l'amplitude.

\begin{attentionbox}
Ne pas confondre :
\begin{itemize}
    \item fréquence \(\rightarrow\) hauteur ;
    \item amplitude \(\rightarrow\) intensité.
\end{itemize}
\end{attentionbox}
\end{corrigebox}

%================================================
\section{Deux problèmes transversaux}
%================================================

%------------------------------------------------
\subsection{Problème transversal 1 -- Du microscopique au macroscopique}
%------------------------------------------------

\begin{enoncebox}
On compare deux liquides A et B de masses molaires voisines.  
Le liquide A est très miscible à l'eau et possède une température d'ébullition relativement élevée.  
Le liquide B est peu miscible à l'eau et s'évapore plus facilement.

On sait que A possède un groupe O--H, alors que B ne possède que des liaisons C--C et C--H.

\begin{enumerate}[label=\arabic*)]
    \item Expliquer pourquoi A et B n'ont pas le même comportement face à l'eau.
    \item Expliquer pourquoi leurs températures d'ébullition peuvent être différentes malgré des masses molaires voisines.
    \item Montrer que la structure moléculaire permet de prévoir des propriétés macroscopiques.
\end{enumerate}
\end{enoncebox}

\begin{corrigebox}
\textbf{1) Miscibilité avec l'eau}

Le liquide A possède un groupe O--H.  
Il peut interagir favorablement avec l'eau, notamment par liaison hydrogène.  
Le liquide B, composé seulement de liaisons C--C et C--H, est globalement apolaire et interagit mal avec l'eau.

\textbf{2) Température d'ébullition}

À masse molaire voisine, la différence ne vient pas principalement de la masse mais de la nature des interactions intermoléculaires :
\begin{itemize}
    \item A : interactions plus fortes, notamment liaisons hydrogène possibles ;
    \item B : interactions plus faibles, essentiellement de Van der Waals.
\end{itemize}
Il faut donc plus d'énergie pour vaporiser A.

\textbf{3) Structure et propriétés}

La structure détermine :
\begin{itemize}
    \item la polarité éventuelle ;
    \item les interactions possibles ;
    \item donc des propriétés comme la miscibilité et la température d'ébullition.
\end{itemize}

\begin{bilanbox}
La physique-chimie devient vraiment intelligible lorsqu'on sait relier :
\[
\text{structure} \rightarrow \text{interactions} \rightarrow \text{propriétés}
\]
\end{bilanbox}
\end{corrigebox}

%------------------------------------------------
\subsection{Problème transversal 2 -- Concevoir une expérience utile}
%------------------------------------------------

\begin{enoncebox}
Un professeur veut faire comprendre à sa classe la différence entre :
\begin{itemize}
    \item une dissolution ;
    \item une dilution ;
    \item une transformation chimique produisant un gaz.
\end{itemize}

Proposer pour chacun des trois cas une expérience simple de laboratoire ou de salle de classe, puis expliquer :
\begin{enumerate}[label=\arabic*)]
    \item ce qu'on observe à l'échelle macroscopique ;
    \item ce qu'on interprète à l'échelle microscopique ;
    \item en quoi les trois situations sont fondamentalement différentes.
\end{enumerate}
\end{enoncebox}

\begin{corrigebox}
\textbf{Exemple de dissolution}

On dissout du sel dans l'eau.
\begin{itemize}
    \item Macroscopiquement : le solide ``disparaît''.
    \item Microscopiquement : les ions se dispersent et sont hydratés.
\end{itemize}

\textbf{Exemple de dilution}

On ajoute de l'eau à une solution colorée de sulfate de cuivre.
\begin{itemize}
    \item Macroscopiquement : la couleur devient moins intense.
    \item Microscopiquement : les espèces dissoutes sont plus dispersées, mais leur quantité dans le prélèvement est inchangée.
\end{itemize}

\textbf{Exemple de transformation chimique avec gaz}

On fait réagir du zinc avec une solution acide.
\begin{itemize}
    \item Macroscopiquement : apparition de bulles.
    \item Microscopiquement : des espèces chimiques disparaissent et de nouvelles espèces apparaissent, notamment \(\mathrm{H_2}\).
\end{itemize}

\textbf{Différence fondamentale}

\begin{itemize}
    \item \textbf{Dissolution} : pas de nouvelle espèce chimique, seulement dispersion d'un soluté.
    \item \textbf{Dilution} : pas de nouvelle espèce chimique, seulement diminution de concentration.
    \item \textbf{Transformation chimique} : apparition et disparition d'espèces chimiques.
\end{itemize}

\begin{bilanbox}
Savoir distinguer ces trois situations est essentiel : elles peuvent parfois se ressembler visuellement, mais leur sens physique est profondément différent.
\end{bilanbox}
\end{corrigebox}

%================================================
\section{Conclusion}
%================================================

\begin{objectifbox}
Pour progresser vraiment en Première spé, il faut s'entraîner à faire cinq choses :
\begin{enumerate}
    \item identifier les grandeurs pertinentes ;
    \item relier les calculs au sens physique ;
    \item distinguer ce qui relève de la structure, des interactions et des propriétés ;
    \item verbaliser clairement un raisonnement ;
    \item vérifier la cohérence d'un résultat.
\end{enumerate}
\end{objectifbox}

\vfill

\begin{center}
\rule{0.82\linewidth}{0.5pt}\\
\textit{Fin du document}
\end{center}

\end{document}